% 2>/dev/null; MSYS_NO_PATHCONV=1 exec docker run --rm -v .:/work -w /work texlive/texlive:latest latexmk -pdf -interaction=nonstopmode -outdir="${0%.tex}" "$0"
% 就活用: Proxima Fusion "Computational Physicist" (Stellarator Design, Munich, L3) の
% 応募フォーム Cover Letter 欄に貼る motivation letter。学会・ジャーナル用ではない。
\documentclass[11pt]{article}

\usepackage[utf8]{inputenc}
\usepackage[T1]{fontenc}
\usepackage[a4paper,margin=2.5cm]{geometry}
\usepackage{hyperref}
\usepackage{parskip}

\pagestyle{empty}

\begin{document}

\begin{flushright}
% TODO: 実名・所在地・電話に差し替える
lzpel \\
misumi3104@gmail.com \\
\url{https://github.com/lzpel} \\
Tokyo, Japan \\
\today
\end{flushright}

\textbf{Application for Computational Physicist --- Stellarator Design, Proxima Fusion (Munich)}

Dear Proxima Fusion team,

I am applying for the Computational Physicist role because it describes, almost
verbatim, the work I have chosen to do on my own time: building simulation tools
from governing equations to validated results, rather than running existing codes.
I am a software engineer, not a fusion physicist by training --- I want to be
upfront about that --- and I believe the fastest way to demonstrate that I can
contribute to Stellaris-class design work is to show the tools I have already
built and the one I am building now.

My current project, \texttt{mhd-tbr-stell}
(\url{https://github.com/lzpel/mhd-tbr-stell}), is an open-source pipeline that
takes a VMEC equilibrium (\texttt{wout.nc}), generates liquid-metal flow channels
conforming to offset magnetic surfaces, and returns the plant-level feasibility
figures of a breeding blanket: the tritium breeding ratio via OpenMC on CAD
geometry, and the MHD pressure drop and pumping-power fraction via a
correlation-based duct network anchored by resolved OpenFOAM (epotFoam) MHD
simulations. The validation plan is deliberately classical --- Hartmann, Shercliff,
and Hunt flows reproduced before any stellarator geometry is touched --- because I
take the view that a fast, lightweight model is only as credible as the
high-fidelity anchor behind it. A preprint is in preparation; the repository
documents the design and the three-month execution plan.

The foundation underneath it is \texttt{cadrum}
(\url{https://github.com/lzpel/cadrum}), a Rust CAD library I wrote and published
on crates.io: a headless, statically linked OpenCASCADE kernel that builds
parametric B-rep solids (sweeps along twisted spines, lofted magnetic-surface-like
bodies, thin-walled shells), reads and writes STEP with per-solid material
colouring, and runs both natively and in the browser via WebAssembly. Shipping it
required owning everything the role's software half implies: numerical geometry,
C++/Rust interop, cross-platform CI, reproducible Docker toolchains, and
documentation good enough for strangers to build on.

On the LLM requirement: agentic LLM tooling is my daily development environment,
and I design for it rather than merely with it --- \texttt{cadrum} ships a
fixed-layout multiview rendering intended specifically as a state snapshot for
LLM-driven design loops. Working this way every day has also taught me its failure
modes, which is why every generated change in my projects lands behind
deterministic verification: analytic benchmarks, reference-file comparisons, and
round-trip tests.

What I would bring to Munich is tool-building velocity with numerical honesty; what
I would ask from Proxima is the plasma-physics depth I do not yet have. I learn
fastest by shipping into review, and I would be glad to walk through any of the
code above in a technical panel.

Thank you for your consideration.

Sincerely, \\
lzpel % TODO: 実名に差し替える

\end{document}
